% Bảng tên cài đặt. Nguồn: include/lock.h (PC_FOR_EACH_LOCK), bench/bench_barrier.cpp (danh sách barrier),
% bench/bench_counter.cpp (PC_COUNTERS_ATOMIC, PC_COUNTERS_LOCK), include/workload.hpp (pc::opaque).
\section{Tên các cài đặt}

Bảng~\ref{tab:impl:names} cho tên của mọi cài đặt, đúng như trong hình, bảng và log. Lock và barrier
thư viện chỉ dùng để so sánh.

\begin{table}[H]\centering\scriptsize\setlength{\tabcolsep}{3pt}
\begin{tabularx}{\linewidth}{@{}l>{\raggedright\arraybackslash}Xl@{}}
\toprule
\textbf{Tên} & \textbf{Mô tả} & \textbf{Hệ điều hành} \\
\midrule
\multicolumn{3}{@{}l}{\emph{Lock}} \\
\PCname{tas} & TAS: lặp \texttt{exchange(acquire)}, nhả bằng \texttt{store(release)} & mọi máy \\
\PCname{ttas} & TTAS: chờ bằng \texttt{load(relaxed)}, rồi mới \texttt{exchange} & mọi máy \\
\PCname{mcs} & MCS: hàng đợi FIFO, mỗi thread chờ trên nút riêng & mọi máy \\
\PCname{tas_bo_{sleep,spin,yield}} & TAS, gọi chính sách chờ sau mỗi lần thử hỏng (\texttt{sleep\_for} nhân đôi; quay ngẫu nhiên; quay 64 lần rồi \texttt{yield}) & mọi máy \\
\PCname{ttas_bo_{sleep,spin,yield}} & TTAS, gọi chính sách chờ khi thua \texttt{exchange} & mọi máy \\
\PCname{std_mutex}, \PCname{pthread_mutex} & \texttt{std::mutex}, \texttt{pthread\_mutex\_t} & mọi máy \\
\PCname{os_unfair_lock} & \texttt{os\_unfair\_lock} của macOS & macOS \\
\PCname{pthread_spinlock} & \texttt{pthread\_spinlock\_t} & Linux \\
\PCname{omp_lock} & \texttt{omp\_lock\_t} (cần bản dựng có OpenMP) & Linux \\
\midrule
\multicolumn{3}{@{}l}{\emph{Barrier}} \\
\PCname{sr} & sense-reversing barrier, chờ bằng quay & mọi máy \\
\PCname{sr_bo_{sleep,spin,yield}} & \PCname{sr}, gọi chính sách chờ như lock giữa hai lần đọc cờ & mọi máy \\
\PCname{sr_block} & \PCname{sr}, chờ chặn bằng \texttt{std::atomic::wait} và \texttt{notify\_all} & mọi máy \\
\PCname{std_barrier} & \texttt{std::barrier} (C++20) & mọi máy \\
\PCname{condvar} & mutex, condition variable và bộ đếm thế hệ & mọi máy \\
\PCname{pthread_barrier} & \texttt{pthread\_barrier\_t} (macOS không có) & Linux \\
\PCname{omp_barrier} & \texttt{\#pragma omp barrier} (cần bản dựng có OpenMP) & Linux \\
\midrule
\multicolumn{3}{@{}l}{\emph{Counter}} \\
\PCname{faa}, \PCname{faa_seqcst} & \texttt{fetch\_add}, thứ tự \texttt{relaxed} hoặc \texttt{seq\_cst} & mọi máy \\
\PCname{cas}, \PCname{cas_seqcst} & vòng CAS không backoff, thứ tự \texttt{relaxed} hoặc \texttt{seq\_cst} & mọi máy \\
\PCname{cas_{sleep,spin,yield}} & vòng CAS, gọi chính sách chờ sau mỗi lần CAS hỏng & mọi máy \\
\PCname{lock_X} & biến \texttt{long} trong critical section của lock \PCname{X} (mọi lock ở trên) & như lock \PCname{X} \\
\bottomrule
\end{tabularx}
\caption{Tên các cài đặt. Hai máy macOS (\PCname{apple-m5}, \PCname{C4M7Y670NT}) có 12 lock, 7 barrier
và 19 counter. \PCname{LAP1184} (Linux, có OpenMP) có 13 lock, 9 barrier và 20 counter.}
\label{tab:impl:names}
\end{table}
% nguồn: results/{apple-m5,C4M7Y670NT}/test.log (12 lock, 7 barrier, 19 counter); results/LAP1184/test.log (13, 9, 20).

Benchmark và bộ kiểm tra cho các giá trị mà compiler có thể gộp hoặc xoá đi qua \texttt{pc::opaque()}:
đối số $N$ của mỗi lời gọi \texttt{is\_prime}, địa chỉ counter ở mỗi vòng của bài kiểm tra counter, và
giá trị đọc ra của biến đếm trong bài kiểm tra lock. Hàm này chỉ là một lệnh \texttt{asm volatile} rỗng
với ràng buộc \texttt{+r}. Compiler phải coi giá trị trong thanh ghi đã bị sửa, nên không chứng minh được
hai lời gọi giống nhau và không gộp hay xoá chúng. Hàm không sinh lệnh máy nào.
% nguồn: include/workload.hpp (opaque); bench/bench_lock.cpp:81, bench/bench_counter.cpp:102,
% tests/test_counter.cpp:326, tests/test_locks.cpp:116.
